\section{Архитектура модели и выгода гибрида}
\label{sec:models}

Сколько даёт гибрид, зависит и от устройства модели. Важны две величины: сколько клиент тратит на чтение порции на
глубине и сколько весит дельта, которую сервер присылает вместо этого. Таблица~\ref{tab:arch} сравнивает четыре
плотные модели одной весовой категории: Qwen3.8-27B, Muse Glimmer 30B \cite{glimmer30b}, Gemma 4 31B
\cite{gemma4} и Qwen3-32B \cite{qwen3}.

\begin{table}[t]
\centering
\caption{Дельта и чтение порции для плотных моделей около 30 миллиардов параметров. Оценки по опубликованным
конфигурациям и формулам приложения~\ref{app:formulas}; кэш в 8 битах. Время чтения на клиенте пересчитано с замера
Qwen3.8-27B на M3 Max (41.2 токена в секунду на глубине 131 тысяча) по числу операций на токен. Экономия на вызове
равна времени чтения на клиенте минус время передачи дельты; у Gemma 4 при 100~Мбит/с передача дольше чтения, и клиент
читает сам. Qwen3-32B без расширения контекста работает до 41 тысячи токенов.}
\label{tab:arch}
\footnotesize
\begin{tabular}{lrrrrr}
\toprule
 & Qwen3.8-27B & Muse Glimmer & Gemma 4 & Qwen3-32B \\
 &  & 30B & 31B & \\
\midrule
Слои полного внимания & 16 из 64 & 13 из 52 & 10 из 60 & 64 из 64 \\
Остальные слои & линейные & окно 2048 & окно 1024 & нет \\
Контекст, тыс. токенов & 262 & 131 & 262 & 41 \\
Кэш на токен, КБ & 34.8 & 7.1 + 21.2 & 21.8 + 435 & 139 \\
Кэш окна или состояние целиком, МБ & 157 & 43 & 446 & 0 \\
Дельта на 1000 токенов, МБ & 192 & 28 & 457 & 139 \\
Передача при 300 Мбит/с, с & 5.1 & 0.8 & 12.2 & 3.7 \\
Чтение 1000 токенов на клиенте, с & 24 & 19 & 34 & 77 \\
\midrule
Экономия на вызове при 100 Мбит/с, с & 8.9 & 17.0 & 0 & 66.0 \\
Экономия на вызове при 300 Мбит/с, с & 19.2 & 18.6 & 21.6 & 73.4 \\
Экономия на вызове при 1 Гбит/с, с & 22.7 & 19.1 & 30.1 & 76.0 \\
\bottomrule
\end{tabular}
\end{table}

\textbf{Чтение на глубине.} С глубиной дорожают только слои полного внимания: каждый новый токен в них смотрит на весь
контекст. Слои линейного внимания и слои со скользящим окном смотрят на фиксированный объём, и их цена от глубины не
зависит. У Qwen3-32B полное внимание во всех 64 слоях, и тысяча токенов на глубине 131 тысяча (с расширением контекста) требует втрое больше вычислений, чем у Qwen3.8-27B; клиент читает её больше минуты.

\textbf{Вес дельты.} Дельта складывается из кэша слоёв полного внимания для новых токенов, кэша окна для новых позиций, но не больше
размера окна, и рекуррентного состояния. У
Qwen3.8-27B постоянную часть составляет рекуррентное состояние: 157~МБ при каждом оффлоаде. У моделей со скользящим окном сервер передаёт кэш окна для новых позиций, и его вес задают число голов ключей и
значений и их размерность. Muse
Glimmer держит в каждом слое окна по две головы размерности 128, и дельта на тысячу токенов весит 28~МБ: канал в
300~Мбит/с передаёт её меньше чем за секунду. Gemma 4 31B держит по 16 голов размерности 256, и та же дельта весит
457~МБ, больше чем у Qwen3.8-27B. Сжатый кэш внимания делает дельту лёгкой и у моделей MoE: у GLM-5.3-Flash он
занимает 6~КБ на токен (раздел~\ref{sec:million}).

\textbf{Какая модель удобнее гибриду.} Ответ зависит от канала и глубины. На канале 100~Мбит/с выигрывает Muse
Glimmer: её дельта лёгкая, и сервер экономит клиенту 17 секунд на порции против 9 у Qwen3.8-27B, а Gemma 4 на таком
канале выгоднее читать самой. На 300~Мбит/с три модели почти равны, 19--22 секунды. На гигабите больше всех экономит
Gemma 4, 30 секунд: её слои полного внимания с большими головами делают чтение на клиенте самым дорогим. Глубже 131
тысячи токенов Glimmer не работает, а у Qwen3.8-27B и Gemma 4 чтение тысячи токенов на глубине 262 тысячи растёт до 36
и 53 секунд, и выигрыш гибрида вместе с ним. Qwen3-32B экономит больше всех на каждой порции, но без расширения
контекста работает только до 41 тысячи токенов. Из четырёх моделей Qwen3.8-27B ближе всего к середине: умеренная
дельта, длинный контекст и выигрыш на любом канале от 100~Мбит/с.

\textbf{Физика выигрыша.} Предел выигрыша со стороны канала задаёт одно отношение: сколько вычислений prefill сервер
снимает с клиента на каждый байт дельты. У клиента своё отношение: сколько он сам успевает посчитать, пока по каналу
идёт один байт. Если первое меньше второго, оффлоад проигрывает идеальному локальному движку, который работает на
пике производительности клиента, даже при мгновенном сервере. Если первое больше, оффлоад может выиграть, и частное
двух отношений даёт верхнюю границу выигрыша: к передаче дельты добавляются чтение на сервере, подъём копии сессии,
вставка дельты и последний токен (приложение~\ref{app:formulas}). Таблица~\ref{tab:intensity} даёт эту границу для
трёх моделей работы. Многое решает постоянная часть дельты: на порции в тысячу токенов у Qwen3.8-27B 157~МБ из 192
приходятся на состояние, и на глубине 110 тысяч граница M5 Max на гигабите близка к единице. На порции в 10 тысяч
граница равна 4.0; с чтением на сервере, но без подъёма копии, вставки и последнего токена, выигрыш против
идеального движка не больше 3.4 раза. На 10~Гбит/с граница вырастает до 40, но главным слагаемым становится чтение на
сервере, и выигрыш не больше 14 раз. Выше 22.5 раза он не поднимется ни на каком канале: во столько раз сервер с B300 в наших оценках
считает prefill быстрее пика M5 Max. Qwen3.8-Flash-Next с 6 миллиардами активных параметров читает дёшево, и на
гигабите граница M5 Max больше единицы только на крупных порциях, 1.5--1.8. Сегодняшние движки на Маке работают
далеко от пика, поэтому на практике оффлоад выигрывает и там, где граница меньше единицы. Такой выигрыш держится на
отставании движков и сократится вместе с ним; выигрыш там, где граница заметно больше единицы, от клиентских
движков не зависит.

\begin{table}[t]
\centering
\caption{Вычисления prefill, которые сервер снимает с клиента на каждый байт дельты, и граница выигрыша со стороны
канала: во сколько раз чтение порции на идеальном локальном движке, который работает на пике производительности
клиента, дольше передачи её дельты, на глубине 110 тысяч и миллион токенов. Меньше единицы: оффлоад проигрывает
идеальному движку при любом сервере. Больше единицы: оффлоад может выиграть, но чтение на сервере, подъём копии
сессии и вставка дельты уменьшают выигрыш (текст). Пик M5 Max 60~TFLOPS, M5 Ultra 140~TFLOPS; сервер с B300
читает со скоростью 1.35~PFLOPS, 30\% пика FP8 (оценки). На 10~Гбит/с граница в десять раз выше. Состояние Qwen3.8-27B в 32 битах, у MoE-моделей в 16. Для Qwen3.8-27B глубина
миллион лежит за пределами штатного контекста. Оценки по формулам приложения~\ref{app:formulas}, кэш в 8 битах.}
\label{tab:intensity}
\footnotesize
\begin{tabular}{lrrr}
\toprule
Модель, порция & МФЛОП на байт & M5 Max, 1~Гбит/с & M5 Ultra, 1~Гбит/с \\
\midrule
Qwen3.8-27B, 1 тыс. & 0.51 / 2.3 & 1.1 / 4.9 & 0.5 / 2.1 \\
Qwen3.8-27B, 10 тыс. & 1.9 / 8.8 & 4.0 / 18 & 1.7 / 7.9 \\
Qwen3.8-Flash-Next, 1 тыс. & 0.19 / 0.23 & 0.4 / 0.5 & 0.2 / 0.2 \\
Qwen3.8-Flash-Next, 10 тыс. & 0.73 / 0.88 & 1.5 / 1.8 & 0.7 / 0.8 \\
GLM-5.3-Flash, 1 тыс. & 0.60 / 0.85 & 1.3 / 1.8 & 0.5 / 0.8 \\
GLM-5.3-Flash, 10 тыс. & 3.6 / 5.1 & 7.6 / 11 & 3.2 / 4.6 \\
\bottomrule
\end{tabular}
\end{table}

\textbf{Идеальная модель для гибрида.} Цену чтения задают слои полного внимания и число голов запросов, вес кэша
задаёт число голов ключей и значений. Гибриду выгодна модель, которая читает глубокий контекст дорого и отдаёт его
дёшево:
\begin{itemize}
\item полное внимание в части слоёв, чтобы сервер снимал с клиента много работы, и сжатый кэш в этих слоях: одна-две
головы ключей и значений или латентный кэш, как MLA;
\item небольшая постоянная часть дельты: рекуррентное состояние или кэш окна на десятки мегабайт, в 16 битах;
\item собственный контекст от 256 тысяч до миллиона токенов, потому что выигрыш растёт с глубиной;
\item размер, при котором модель с кэшем нескольких сессий помещается в память клиента и быстро генерирует;
\item детерминированные ядра вычислений, чтобы клиент мог сверять дельты точно (раздел~\ref{sec:security}).
\end{itemize}
Например, если бы у Qwen3.8-27B в слоях полного внимания был кэш размером с кэш GLM-5.3-Flash, около 6~КБ на
токен, а состояние передавалось в 16 битах, дельта на тысячу токенов весила бы около 84~МБ вместо 192. На канале
100~Мбит/с сервер тогда экономил бы клиенту около 18 секунд на порции вместо 9. Это оценка по тем же формулам.

\textbf{Вывод.} Выбирая модель для гибрида, достаточно двух чисел: сколько стоит чтение тысячи токенов на глубине и
сколько весит дельта на тысячу токенов. Первое задаёт выигрыш, второе задаёт канал, который нужен, чтобы этот выигрыш
получить.
